\subsection{All terminal triples (1,a,b) with a at least four}\label{proof:Q5-C173}
% Common local aliases for the former repeated C168.1--C168.3 prefaces.
Use the original-measure calculus of
Section~\ref{sec:odd-common-calculus} and the complete root-pair
certificates of Section~\ref{sec:odd-common:fourth}.
For a root \((1,2,v)\), write \(m=v-2\), \(G_m=E+zD\),
\(\beta=\beta_{v-1}=\int D\,d\nu\),
\(\beta_v=\int E\,d\nu\),
\(P_+=h_2f_v\), \(P_-=\bar f_2W_v\),
\(\ell_+=\ell_a\), and \(\ell_-=\ell_s\).
The two vectors, their signed means and full constants are
\eqref{eq:odd-common:root-standard}--\eqref{eq:odd-common:root-losses}.
At \(v=3\), \((E,D)=(1,0)\) and \(\beta=0\).
All subsequent direction choices and packet estimates retain their
stated integer and parameter ranges.
Here the root coordinates are represented by $G_{v-2}=E_0+zD_0$. The ordinary-fourth bridge of Q5-C084, Section 2, has cost $11d$ for every positive outer length. Its terminal $(1,a)$ row, $a\ge4$, uses the identical $P_{v,L}$ below and has rate $705/128$. This is the actual same-direction high-domain comparison.
\paragraph{C173.0 Fixed original object and direction}

Every arm length and multiplicity, subdivision vertex, independent positive core length, original edge count and original joint label law is retained.

% Replace a pure model-definition block only; keep local root/terminal data.
Use the original finite-space model and arm messages of
Section~\ref{sec:odd-common-calculus}.
When the original \(K\)-notation is used, \(K=-B\) and
\(g=K\one=-u\); in either notation
\(\delta_j=\|B^j\one\|_2^2=\|K^j\one\|_2^2\).
All means, norms, adjoints and spectral integrals use the same
original probability measure.
The root residual is \((1,\allowbreak \allowbreak  2,\allowbreak \allowbreak  v)\), \(v\ge3\), and the terminal residual \((1,\allowbreak \allowbreak  a,\allowbreak \allowbreak  b)\), \(4\le a<b\). Original positive \(L,\allowbreak \allowbreak  M\) are independent. Always use the same actual terminal and the integer-selected root pair
\[
F_{ab}=qf_a h_b,\qquad
P_{v,L}=\begin{cases}h_2f_v,&L=1,\ v\ge4,\\
\bar f_2W_v,&\text{otherwise},\end{cases}
\qquad H=pP_{v,L}T^LF_{ab}.                         \tag{C173.T183.1}
\]

For all listed original integers and legal kernels,
\[
\boxed{\begin{gathered}
EH^4\le e^{-705d/128},\qquad\\
I_q(T^MH)^3\le e^{-d/64},\quad\\
I_q(Y)=E[q^{-1/3}|Y|^{8/3}].
\end{gathered}}\tag{C173.T183.2}
\]
Genuine equal-pair decorations and simultaneous reversal are included. The moment is the ordinary fourth moment with the actual unit arm \(p\) and complete pair \(P\).

The four arm types have their fixed directions. Root and terminal unit arms give \(p,\allowbreak \allowbreak  q\); inner edges point terminal to root, outer edges root to receiver. All choices precede \(\mu,\allowbreak \allowbreak  K,\allowbreak \allowbreak  d\).

The entire low domain \(0<d\le1/4\) has common rate at least \(349/60\), exceeding \(705/128\) by \(593/1920\). The high domain compares the same direction with the Q5-C084 \(qf_a\) row.


\paragraph{C173.2 Four actual beta constants and every terminal pair}

Through the low-domain payment take \(0<d\le1/4\). Let \(n_0=b-a\ge1\), and define
\[
c=\beta_{b-1}-\beta_{a-1},\quad
\Gamma=\beta_a+\beta_{a-1}+\beta_b+\beta_{b-1},
\]
\[
R=r(B)u,\qquad r(z)=-c+(-1)^a z^aG_{n_0}(-z).       \tag{C173.T183.7}
\]
Retain \(\beta_af_0,\allowbreak \allowbreak  \beta_{a-1}f_1\) in the finite convolution for \(f_a\); retain \(\beta_{b-1}f_0\) for \(f_{b-1}\) and apply positive \(T\). With \(K=-B\), \(S_n=\sum_{j=0}^nK^j\one\),
\[
f_a\le S_a-\beta_a-\beta_{a-1}p,\qquad
h_b\le TS_{b-1}-\beta_{b-1}q,
\]
\[
TS_{b-1}=1-\beta_b-\sum_{j=1}^{b}K^j1.
\]
Adding gives the exact envelope
\[
f_a+h_b\le2-\Gamma+R,\qquad F_{ab}\le q e^{R-\Gamma}. \tag{C173.T183.8}
\]
AM--GM acts on the two actual nonnegative arms; their upper sum remains nonnegative, and \((1+s/2)^2\le e^s\).

The beta difference is a finite decreasing alternating sum beginning at index \(\lfloor a/2\rfloor\), so
\[
|c|\le \delta_{\lfloor a/2\rfloor}.
\]
For \(a=4,\allowbreak \allowbreak  5\), \(\beta_{a-1}=d\), hence \(c\le0,\allowbreak \allowbreak |c|\le \delta_2\); for \(a\ge6\), \(|c|\le \delta_3\). The corresponding scalar lower bounds are
\[
\Gamma\ge
\begin{cases}(51/16)d,&a=4,\\
(89/32)d,&a=5,\\
(19/8)d,&a\ge6.
\end{cases}                                        \tag{C173.T183.9}
\]
They retain respectively \(2d+2\tau,\allowbreak \allowbreak  d+3\tau,\allowbreak \allowbreak  4\tau\), preserving the extra short-arm loss.

The original row bound gives
\[
\begin{gathered}
|R|\le\theta d+(6/5)\theta^{3/2}\sqrt d\\
<13/128+507/3200=13/50<1/2.
\end{gathered}\tag{C173.T183.10}
\]
Use \(\sqrt\theta<13/20\). On \(|t|\le1/2\), \(e^t\le1+t+3t^2/5\), since its integral remainder coefficient increases in \(t\), and \(e^{1/2}\le79/48+5/1728=2849/1728<33/20\). Keeping the actual mean, \(q\le2\), the kernel bound two and row Cauchy give
\[
\begin{gathered}
T(qe^R)\le \bar f_2+BR+\delta,\qquad\\
\delta=ER+E(uR)+\sqrt{dER^2}+(12/5)ER^2.
\end{gathered}\tag{C173.T183.11}
\]
Expand \(T(qR)=ER+E(uR)+BR+B(uR)\); bound only the final term by \(\sqrt{dER^2}\). The quadratic remainder is at most \((12/5)ER^2\).

\paragraph{C173.3 Three first-row bounds on the full low domain}

Write \(G_{b-a}(-z)=U-zV\), with \(0\le U,\allowbreak \allowbreak  V\le1\) for every \(n_0\ge1\), including the actual \(V=0\) at one. Write \(r=A_0+zC_0\), and let \(\rho\) bound the spectral modulus of this finite polynomial.

\paragraph{C173.3.1 The case \(a=4\)}

\[
A_0=-c+x^2U,\quad C_0=-x^2V,\quad
\rho\le\theta/4+(6/5)\theta^2=767/2560<3/10.
\]
The actual means obey \(ER\le d3\) and \(E(uR)\le |c|d+d3,\allowbreak \) with \(ER^2\le \rho^2\) d. Thus
\[
\begin{gathered}
\delta/d\le2\theta^2+\theta/4+3/10+27/125\\
<1/3+5/48+3/10+27/125=1907/2000<1.
\end{gathered}\tag{C173.T183.12}
\]
Take \(D=d\). Also retain the exact signs \(D_1\le0,\allowbreak \allowbreak  D_2\le0\).

\paragraph{C173.3.2 The case \(a=5\)}

\[
A_0=-c-x^3V,\quad C_0=-x^2U.
\]
\[
\rho\le\theta/4+(6/5)\theta^{5/2}
<5/48+13/100=281/1200<1/4.
\]
Now \(ER\le d3\) and \(E(uR)\le |c|d,\allowbreak \) whence
\[
\begin{gathered}
\delta/d\le\theta^2+\theta/4+1/4+3/20\\
=273/1024+2/5<2/3.
\end{gathered}\tag{C173.T183.13}
\]
The final margin is \(1/15360\). Take \(D=2d/3\), retaining \(D_2\le0,\allowbreak \allowbreak  D_1\le \delta_4\).

\paragraph{C173.3.3 Every a at least six, both parities and every finite gap}

For \(a=2h,\allowbreak \) h at least three,
\[
A_0=-c+(-1)^h x^hU,\quad C_0=-(-1)^h x^hV.
\]
For \(a=2h+1,\allowbreak \) h at least three,
\[
A_0=-c+(-1)^{h+1}x^{h+1}V,\quad
C_0=(-1)^{h+1}x^hU.
\]
Thus \(|A_0|\le|c|+x^3\), \(|C_0|\le x^3\), and
\[
\rho\le\theta^2/4+(6/5)\theta^3=9971/81920<1/8.
\]
The actual mean sum is at most \(\delta_3d+2\delta_4\). Using \(\theta^2<1/6\), \(\theta^3<7/100\),
\[
\begin{gathered}
\delta/d\le\theta^2/4+2\theta^3+1/8+3/80\\
<1/24+7/50+13/80=413/1200<7/20.
\end{gathered}\tag{C173.T183.14}
\]
Take \(D=7d/20\). This covers all integer lengths and the whole low domain with the same directions.

\paragraph{C173.4 Every finite signed mean and gamma sign}

For i at least one define \(D_i=EB^iR.\) The exact signs are
\[
\begin{gathered}
D_{2j-1}=(-1)^j\int x^{j-1}A_0,\qquad\\
D_{2j}=(-1)^{j+1}\int x^jC_0.
\end{gathered}\tag{C173.T183.15}
\]
In particular \(D_1=-\int A_0\), \(D_2=\int xC_0\). Uniformly in \(a\ge4\), \(|A_0|\le \delta_2+x^2\), \(|C_0|\le x^2\), so
\[
|D_1|\le \delta_2d+\delta_3
\le(\theta/4+\theta^2)d=(273/1024)d<d/2,
\]
\[
|D_2|\le \delta_4\le\theta^3d<d/12,
\]
\[
|D_i|\le \delta_2^2+\delta_4<d/8\quad(i\ge3).                \tag{C173.T183.16}
\]
For odd \(i\ge3\), the first bound is \(|c|\delta_2+\delta_4\); for even \(i\ge4\), it starts at \(\delta_5\). All later moments decrease since \(x\le\theta<1\). Use \(\theta^2/4<1/24\), \(\theta^3<1/12\). Put \(\gamma_j=\beta_j-D_{j-2}\); since \(\beta_3=\beta_4=d\),
\[
\begin{gathered}
\gamma_3\ge d/2,\quad\gamma_4\ge11d/12,\quad\\
\boxed{\gamma_j\ge15d/32\quad(j\ge3).}
\end{gathered}\tag{C173.T183.17}
\]
For \(j\ge5\), use \(\kappa-1/8=15/32\). Thus every discarded path term has its proved nonnegative coefficient.

Keep these stronger short-length signs:
\[
\begin{array}{c|cc}
 &4D_1&4(D_1+D_2)\\ \hline
a=4&\le0&\le0\\
a=5&<d/3&<d/3\\
a\ge6&<d/2&<3d/5
\end{array}                                        \tag{C173.T183.18}
\]
For \(a=4\), use \(A_0\ge0,\allowbreak \allowbreak  C_0\le0\). For \(a=5\), \(D_1=cd+\int x^3V\le \delta_4\), \(D_2\le0\), and \(4\theta^3<1/3\). For \(a\ge6\),
\[
|D_1|/d\le\theta^2/4+\theta^3<67/600,\quad
|D_2|/d\le\theta^4<1/36,
\]
Hence \(4|D_1|<67d/150<d/2\), and \(4(|D_1|+|D_2|)<251d/450<3d/5\), using \(\theta^4<1/36\).

\paragraph{C173.5 Actual positive propagation and the full root payment}

Finite induction using the actual positive T gives
\[
\begin{gathered}
T^LF_{ab}\le e^{-\Gamma}Y_L,\quad\\
Y_L=\bar f_{L+1}+B^LR+D\bar f_{L-1}\\
+\sum_{i=1}^{L-1}D_i \bar f_{L-1-i}.
\end{gathered}\tag{C173.T183.19}
\]
This follows step by step from \(T(B^iR)=D_i+B^{i+1}R\). The right side dominates a positive actual message. For \(L\ge2\), exact convolution gives
\[
Y_L=1+G_{L+1}(B)u+B^LR+D\bar f_{L-1}
 -\sum_{j=3}^{L+1}\gamma_j\bar f_{L+1-j}.
\]
Retain the final original \(\bar f_0\) and use \(\bar f_{L-1}\le2\):
\[
\begin{gathered}
Y_L\le1+u+2D-\gamma_{L+1}+BQ_L(B)u,\quad\\
Q_L=G_L+z^{L-1}r,\quad E[BQ_L(B)u]=-\gamma_{L+2}.
\end{gathered}\tag{C173.T183.20}
\]
At \(L=1\), use exactly \(Y_1=1+u+D+B(u+R)\), with \(D\) once.


Use the complete root envelopes of Lemma~\ref{lem:odd-common:root-pairs} and the actual $p$ curvature in \eqref{eq:odd-common:curvature}. For L at least two the full joint vector gives
\[
\begin{gathered}
\log EH^4\le-4\Gamma+8D\\
-4(\gamma_{L+1}+\gamma_{L+2})-4\ell_-\\
+(17/4)\|[Q_L-zG_n(z)]u\|^2+(10/3)\beta^2.
\end{gathered}\tag{C173.T183.24}
\]
At \(L=1\) the exact corresponding expression is
\[
\begin{gathered}
\log EH^4\le-4\Gamma+4D-4d+4D_1-4\ell_\sigma\\
+(17/4)\|[1+r-zG_n(\mp z)]u\|^2+(10/3)\beta^2.
\end{gathered}\tag{C173.T183.25}
\]
The positive root uses \(G_n(-z)\), only when \(v\ge4\); the negative-root first row occurs only at \(v=3\). At \(L=2\), retain \(\gamma_3+\gamma_4=2d-D_1-D_2\). The curvature cost is at most \(5d/6\), and exactly zero at \(v=3\).

\paragraph{C173.6 Three root bases on the entire finite-frequency interval}

Temporarily omit \(r\), and let \(\ell_\sigma(x)\) be the same-frequency root-loss integrand. Each inequality below holds on \([0,\allowbreak \allowbreak  \theta]\), before integration against the original \(\nu\).

\paragraph{C173.6.1 The first connector in its actual two root directions}

At \(v=3\) the negative-root base is 1-z, with loss integrand one and beta zero. Thus
\[
\begin{gathered}
-4\ell_-+(17/4)|1-z|^2\le253/128<2,\qquad\\
|1-z|<6/5.
\end{gathered}\tag{C173.T183.26}
\]
For \(v\ge4\), the positive-root base is \(1-zG_n(-z)=1-xD_0-zE_0\). Since \(0\le1-xD_0\), \(E_0\le1,\allowbreak \allowbreak  D_0\ge1-x\),
\[
|1-zG_n(-z)|^2\le(1-x+x^2)^2+x\le1.
\]
The difference divided by \(x\) is \(-1+3x-2x^2+x^3\), increasing with derivative \(3-4x+3x^2>0\); its endpoint is \(-1451/32768<0\). Also \(\ell_+\ge(1-x)(2-x)\ge969/1024\), hence
\[
-4\ell_++(17/4)|1-zG_n(-z)|^2\le119/256<1/2.        \tag{C173.T183.27}
\]
Only the actual positive-root range n at least two is used.

\paragraph{C173.6.2 The second connector with the negative root}

The base is \(G_2-zG_n=1+xD_0+z(1-E_0)\). Its joint derivative in \(D_0\) is at most \((1+x)(17x/2-4)<0\), and in \(E_0\) at most \(-4\). For \(n\ge2\), set both to \(1-x\), obtaining
\[
\begin{gathered}
-15/4+(25/2)x-x^2/4-(17/4)x^3+(17/4)x^4\\
\le-15/4+(25/2)\theta=85/64<3/2.
\end{gathered}\tag{C173.T183.28}
\]
At \(n=1\), the actual base is one and the joint expression \(1/4\). Uniformly in \(n\),
\[
\begin{gathered}
|G2-zG_n|^2\le(1+x)^2+x^3\\
\le66997/32768<(23/16)^2.
\end{gathered}\tag{C173.T183.29}
\]

\paragraph{C173.6.3 Every finite length at least three with the negative root}

Write \(G_L=U_L+zV_L.\) Its joint base is
\[
-4[E_0+(1+x)D_0]
+(17/4)[(U_L+xD_0)^2+x(V_L-E_0)^2].
\]
The \(D_0\)-derivative is negative as above, and the \(E_0\)-derivative at most \(-4+17x^2/2<0\). At \(n\ge2\), set both to \(1-x\), use \(U_L\le1-x+x^2\), \(|V_L-(1-x)|\le x\), giving
\[
-15/4+4x+4x^2+(17/4)x^3\le-7/32<1/4.
\]
At \(n=1\), the square is at most \((1-x+x^2)^2+x^3\le1\), because \(2-3x+x^2-x^3>0\); the joint bound is again \(1/4\). The norm coefficient \(23/16\) above holds for every finite \(n,\allowbreak \allowbreak  L\ge3\), retaining \(G_1,\allowbreak \allowbreak  G_2\).

\paragraph{C173.7 Three complete perturbation payments}

At the same frequency use \(|B+\varepsilon|^2\le|B|^2+2|B||\varepsilon|+|\varepsilon|^2\). Take \(\rho=3/10,\allowbreak \allowbreak  1/4,\allowbreak \allowbreak  1/8\) for \(a=4,\allowbreak \allowbreak  5,\allowbreak \allowbreak  \ge6\). The following are costs after multiplication by \(17/4\), per unit \(d\).

The first-row base has modulus below \(6/5\) for the negative root and at most one for the positive root:
\[
\begin{array}{c|cc}
&L1\text{ negative root}&L1\text{ positive root}\\ \hline
a4&1377/400<7/2&1173/400<3\\
a5&901/320<3&153/64<5/2\\
a\ge6&1717/1280<7/5&289/256<6/5
\end{array}                                        \tag{C173.T183.30}
\]
The formulas are \((17/4)[(12/5)\rho+\rho^2]\) and \((17/4)[2\rho+\rho^2]\), respectively.

At \(L=2\), the perturbation is \(zr\), of modulus at most \(\sqrt\theta\rho<(13/20)\rho\). With base bound \(23/16\),
\[
\begin{array}{c|cc}
&|\epsilon|\text{ upper bound}&(17/4)[(23/8)|\epsilon|+|\epsilon|^2]\\ \hline
a4&39/200&203541/80000<13/5\\
a5&13/80&53703/25600<21/10\\
a\ge6&13/160&104533/102400<21/20
\end{array}                                        \tag{C173.T183.31}
\]
The second-row rational margin is \(57/25600\).

For every finite \(L\ge3\), the perturbation \(z^{L-1}r\) has modulus at most \(\rho x\). After integration its full cost is
\[
(17/4)[(23/8)\rho\theta+\rho^2\theta^2]d
\le\begin{cases}
(635817/409600)d<(8/5)d,&a4,\\
(84201/65536)d<(13/10)d,&a5,\\
(165529/262144)d<(13/20)d,&a\ge6.
\end{cases}                                        \tag{C173.T183.32}
\]
The three comparisons are \(635817<655360\), \(842010<851968\), and \(3310580<3407872\).

\paragraph{C173.8 All short connectors and every finite tail}

Substitute the actual \(\Gamma,\allowbreak \allowbreak  D,\allowbreak \allowbreak  D_i\), root beta and full joint base into the two common payment formulas. The table gives positive fourth losses \(\lambda\), with \(\log EH^4\le-\lambda d\), on \(d\le1/4\).


\begingroup\small
\setlength{\tabcolsep}{3pt}
\setlength{\LTleft}{0pt}\setlength{\LTright}{\fill}
\begin{longtable}{@{}p{0.1200\linewidth}p{0.8000\linewidth}@{}}
\caption{All terminal triples (1,a,b) with a at least four: parameter values.}\label{tab:parameters:q5-c173:1}\\
\toprule
Entry & \(L=1,\allowbreak  v=3,\allowbreak \) negative root \\
\midrule\endfirsthead
\toprule
Entry & \(L=1,\allowbreak  v=3,\allowbreak \) negative root \\
\midrule\endhead
\bottomrule\endfoot
\(a=4\) & \(29/4\); \(L=1,\allowbreak  v\ge 4,\allowbreak \) positive root: \(101/12\); \(L=2,\allowbreak \) all v: \(469/60\); \(L\ge 3,\allowbreak \) all v: \(349/60\). \\
\addlinespace[2pt]
\(a=5\) & \(57/8\); \(L=1,\allowbreak  v\ge 4,\allowbreak \) positive root: \(199/24\); \(L=2,\allowbreak \) all v: \(361/40\); \(L\ge 3,\allowbreak \) all v: \(859/120\). \\
\addlinespace[2pt]
\(a\ge 6\) & \(41/5\); \(L=1,\allowbreak  v\ge 4,\allowbreak \) positive root: \(136/15\); \(L=2,\allowbreak \) all v: \(643/60\); \(L\ge 3,\allowbreak \) all v: \(523/60\). \\
\end{longtable}
\endgroup


The following expressions give every rational entry.

For \(a=4,\allowbreak \allowbreak  L=1\), \(4D-4d=0,\allowbreak \allowbreak  4D_1\le0\), and the negative root has \(\beta=0\):
\[
51/4-2-7/2=29/4,\qquad
51/4-1/2-3-5/6=101/12.
\]
At \(L=2\), \(8D-8d=0\), \(4(D_1+D_2)\le0\):
\[
51/4-3/2-13/5-5/6=469/60.
\]
For \(L\ge3\), the adjacent gamma sum is at least \(15d/16\):
\[
51/4-8+15/4-1/4-8/5-5/6=349/60.
\]

For \(a=5,\allowbreak \allowbreak  L=1\), \(4d-4D=4d/3\), \(4D_1<d/3\):
\[
89/8+4/3-1/3-2-3=57/8,
\]
\[
89/8+4/3-1/3-1/2-5/2-5/6=199/24.
\]
At \(L=2\), retain \(4(D_1+D_2)<d/3\):
\[
89/8+8/3-1/3-3/2-21/10-5/6=361/40.
\]
For longer L:
\[
89/8-16/3+15/4-1/4-13/10-5/6=859/120.
\]

For \(a\ge6,\allowbreak \allowbreak  L=1\), \(4d-4D=13d/5\):
\[
19/2+13/5-1/2-2-7/5=41/5,
\]
\[
19/2+13/5-1/2-1/2-6/5-5/6=136/15.
\]
At \(L=2\) and at longer lengths the payments are respectively
\[
19/2+26/5-3/5-3/2-21/20-5/6=643/60,
\]
\[
19/2-14/5+15/4-1/4-13/20-5/6=523/60.
\]

Thus the entire low domain and every original integer satisfy
\[
\boxed{EH^4\le e^{-349d/60}<e^{-705d/128}.}          \tag{C173.T183.33}
\]
The difference from \(705/128\) is \(593/1920\).

\paragraph{C173.9 The same-direction high domain and every outer length}

For \(1/4\le d\le1/3\), \(b\ge5\), so the exact original path-mean bound gives
\[
0\le h_b\le2Ef_{b-1}\le2(1-2d)\le1.
\]
With the same actual F and P, positivity yields
\[
0\le H\le pP_{v,L}T^L(qf_a).
\]
The right side is exactly the Q5-C084 same-direction row for \(a\ge3\), giving rate \(705/128\). Both domains include \(d=1/4\). At zero, leaf integration makes every message one, with no requirement that \(B=0\). The ordinary fourth bound therefore holds on the full original host domain.

For every independent positive outer length, the ordinary fourth bridge gives
\[
\begin{gathered}
I_q(T^MH)^3\le\\
\exp\{[11-2(705/128)]d\}=e^{-d/64}.
\end{gathered}\tag{C173.T183.34}
\]
Its input is the actual ordinary fourth moment.
